\documentclass[12pt,a4paper]{article}

\usepackage[margin=1in]{geometry}
\usepackage[T1]{fontenc}
\usepackage[utf8]{inputenc}
\usepackage{lmodern}
\usepackage{setspace}
\usepackage{hyperref}
\usepackage{booktabs}
\usepackage{longtable}
\usepackage{array}
\usepackage{tabularx}
\usepackage{graphicx}
\usepackage[table]{xcolor}
\usepackage{enumitem}
\usepackage{titlesec}
\usepackage{listings}

\setstretch{1.2}
\setlist[itemize]{noitemsep, topsep=4pt}
\setlist[enumerate]{noitemsep, topsep=4pt}

\hypersetup{
  colorlinks=true,
  linkcolor=blue,
  urlcolor=blue,
  pdftitle={Customer Shopping Behavior Analysis},
  pdfauthor={Data Analyst Portfolio Project}
}

\titleformat{\section}{\large\bfseries}{\thesection.}{0.5em}{}
\titleformat{\subsection}{\normalsize\bfseries}{\thesubsection.}{0.5em}{}

\lstset{
  basicstyle=\ttfamily\small,
  breaklines=true,
  frame=single,
  backgroundcolor=\color{gray!8},
  columns=fullflexible
}

\newcommand{\findingbox}[3]{
\begin{center}
\setlength{\arrayrulewidth}{0.6pt}
\renewcommand{\arraystretch}{1.25}
\begin{tabularx}{0.97\linewidth}{|>{\raggedright\arraybackslash\bfseries}p{0.2\linewidth}|X|}
\hline
{\cellcolor{gray!12}Finding} & #1 \\
\hline
{\cellcolor{gray!12}Business implication} & #2 \\
\hline
{\cellcolor{gray!12}Recommended decision} & #3 \\
\hline
\end{tabularx}
\end{center}
}

\title{\textbf{Customer Shopping Behavior Analysis}\\[0.4em]
\large \textcolor{gray}{Python \textbar{} SQL \textbar{} Power BI}}
\author{}
\date{}

\begin{document}

\maketitle
\thispagestyle{empty}

\begin{abstract}
This report analyzes 3,900 customer shopping records using Python, SQL, and Power BI to evaluate revenue concentration, customer segments, subscription effectiveness, discount behavior, and product performance.
\end{abstract}

\clearpage
\tableofcontents
\clearpage

\section{Executive Context}
Retail teams need fast answers to a few high-value questions: which categories drive revenue, which customer groups matter most, whether retention programs are working, and where promotions may be hurting margin. This project addresses those questions across three tools:
\begin{itemize}
  \item \textbf{Python} for data preparation, quality checks, and feature engineering.
  \item \textbf{SQL} for structured analysis of business questions.
  \item \textbf{Power BI} for stakeholder-friendly storytelling and interactive exploration.
\end{itemize}

\section{Business Objective}
The objective is to identify how customer behavior, product mix, and commercial tactics affect revenue and customer value, then translate those findings into practical decisions for marketing, merchandising, and retention.

The analysis was scoped into five questions:
\begin{enumerate}
  \item Measure the overall size and shape of the business using customer count, revenue, and average purchase metrics.
  \item Compare revenue and spending behavior across customer groups such as gender, age group, and subscription status.
  \item Identify which product categories and individual products contribute most to sales performance.
  \item Test whether discounts, subscriptions, and repeat purchase behavior are associated with higher customer value.
  \item Surface findings in a Power BI dashboard that a non-technical stakeholder can use interactively.
\end{enumerate}

\section{Dataset Overview}
The dataset contains \textbf{3,900 rows} and \textbf{18 columns}, where each row represents a customer purchase record.

\subsection{Core Fields}
The dataset includes the following important dimensions:
\begin{itemize}
  \item \textbf{Customer information:} customer ID, age, gender, location, subscription status.
  \item \textbf{Purchase information:} item purchased, category, purchase amount in USD, season, size, and color.
  \item \textbf{Behavioral information:} previous purchases, frequency of purchases, review rating.
  \item \textbf{Commercial levers:} discount applied, promo code used, shipping type, payment method.
\end{itemize}

\subsection{Initial Data Quality Observations}
Before any analysis, I verified the main quality characteristics of the dataset:
\begin{itemize}
  \item \textbf{Missing values:} 37 missing values were present in \texttt{Review Rating}.
  \item \textbf{Duplicate rows:} 0 duplicate rows were found.
  \item \textbf{Customer uniqueness:} 3,900 unique customer IDs were present, matching the row count.
  \item \textbf{Age range:} customers ranged from 18 to 70 years old.
  \item \textbf{Purchase amount range:} transaction values ranged from 20 to 100 USD.
\end{itemize}

\subsection{High-Level Summary Statistics}
\begin{center}
\begin{tabular}{lr}
\toprule
\textbf{Metric} & \textbf{Value} \\
\midrule
Number of customers / records & 3,900 \\
Total revenue & 233,081 USD \\
Average purchase amount & 59.76 USD \\
Average review rating & 3.75 / 5 \\
Median purchase amount & 60.00 USD \\
Purchase amount standard deviation & 23.69 USD \\
\bottomrule
\end{tabular}
\end{center}

\section{Analytical Workflow}
The project followed a clear end-to-end sequence:
\begin{enumerate}
  \item Load and inspect the CSV in Python.
  \item Identify missing values and apply targeted imputation.
  \item Standardize field names and engineer new analysis features.
  \item Remove redundant information before database loading.
  \item Push the cleaned dataset into MySQL using SQLAlchemy.
  \item Use SQL to answer predefined business questions.
  \item Build a Power BI dashboard that summarizes the most important patterns.
\end{enumerate}

\section{Data Preparation in Python}
\subsection{Loading and Inspecting the Data}
The analysis begins in Python with \texttt{pandas}, using structure review, null checks, and summary statistics.

\subsection{Handling Missing Values}
The only incomplete field was \texttt{Review Rating}, with 37 missing observations. Instead of dropping those rows or filling with a global average, I used category-level median imputation.
\begin{itemize}
  \item review patterns may differ by category;
  \item median imputation is less sensitive to outliers than mean imputation;
  \item keeping the rows preserves sales information for later aggregation.
\end{itemize}

\subsection{Column Standardization}
I standardized column names to snake case and removed spaces and parentheses. For example, \texttt{Purchase Amount (USD)} became \texttt{purchase\_amount\_usd}.

\subsection{Feature Engineering}
Two derived features were added:
\begin{itemize}
  \item \textbf{Age group:} customer age was binned into four quantile-based groups: Young Adult, Adult, Middle-aged, and Senior.
  \item \textbf{Purchase frequency days:} purchase frequency labels such as Weekly, Monthly, and Quarterly were mapped to approximate numeric day intervals.
\end{itemize}

The \texttt{age\_group} field improves segmentation, while \texttt{purchase\_frequency\_days} provides a base for future retention or customer lifetime value analysis.

\subsection{Redundancy Check}
I checked whether \texttt{discount\_applied} and \texttt{promo\_code\_used} carried distinct information. In this dataset, the two fields matched perfectly, so \texttt{promo\_code\_used} was dropped from the cleaned analysis table.

\subsection{Database Loading}
After cleaning, the DataFrame was written into a MySQL database using \texttt{SQLAlchemy} and \texttt{PyMySQL}. Although an earlier project summary mentioned PostgreSQL, the actual implementation in the notebook and SQL script uses MySQL-compatible syntax, so this report follows the real project pipeline.

The cleaned table was stored as \texttt{customer} and used for SQL analysis.

\section{SQL Analysis Design}
The SQL phase translates business questions into structured queries covering revenue, product performance, customer segmentation, and subscription behavior.

\subsection{Business Questions}
\begin{enumerate}
  \item What is the total revenue generated by male vs. female customers?
  \item Which customers used a discount but still spent more than the average purchase amount?
  \item Which are the top 5 products with the highest average review rating?
  \item How do average purchase amounts compare between Standard and Express shipping?
  \item Do subscribed customers spend more than non-subscribed customers?
  \item Which products have the highest percentage of discounted purchases?
  \item How can customers be segmented into New, Returning, and Loyal groups?
  \item What are the top 3 most purchased products within each category?
  \item Are repeat buyers more likely to subscribe?
  \item What is the revenue contribution of each age group?
\end{enumerate}

\subsection{Analytical Coverage}
These questions cover four decision areas:
\begin{itemize}
  \item They compare \textbf{customer value} across groups.
  \item They evaluate \textbf{commercial levers} such as discounts and subscriptions.
  \item They surface \textbf{assortment performance} at both category and item level.
  \item They introduce \textbf{segmentation logic} rather than only descriptive counts.
\end{itemize}

\section{Key Findings and Business Recommendations}
\subsection{Overall Performance}
The business generated \textbf{233,081 USD} in revenue across 3,900 transactions, with an average purchase amount of \textbf{59.76 USD}. The distribution is fairly centered around 60 USD, with the median at 60 USD and a standard deviation of 23.69 USD. This suggests moderate dispersion but no extreme skew.

\subsection{Category Performance}
Category analysis shows a strong concentration of revenue:
\begin{center}
\begin{tabular}{lrrr}
\toprule
\textbf{Category} & \textbf{Revenue} & \textbf{Transactions} & \textbf{Avg. Purchase} \\
\midrule
Clothing & 104,264 & 1,737 & 60.03 \\
Accessories & 74,200 & 1,240 & 59.84 \\
Footwear & 36,093 & 599 & 60.26 \\
Outerwear & 18,524 & 324 & 57.17 \\
\bottomrule
\end{tabular}
\end{center}

Clothing is the dominant category by both revenue and transaction count, contributing nearly half of total revenue. Accessories also performs strongly, while Footwear shows the highest average purchase value among the four categories. Outerwear lags both in volume and revenue, which suggests lower demand or a narrower assortment.

\findingbox
{Clothing and Accessories are the main revenue drivers, while Outerwear contributes the least revenue and volume.}
{Revenue growth is likely to respond fastest to improvements in the strongest categories rather than broad, unfocused optimization across the full catalog.}
{Prioritize Clothing and Accessories in campaign planning, assortment optimization, inventory forecasting, and homepage exposure because improving already strong categories is likely to produce the fastest commercial impact.}

\subsection{Age Group Performance}
Revenue contribution by age group is relatively balanced, but not equal:
\begin{center}
\begin{tabular}{lrrr}
\toprule
\textbf{Age Group} & \textbf{Revenue} & \textbf{Customers} & \textbf{Avg. Purchase} \\
\midrule
Young Adult & 62,143 & 1,028 & 60.45 \\
Middle-aged & 59,197 & 986 & 60.04 \\
Adult & 55,978 & 942 & 59.42 \\
Senior & 55,763 & 944 & 59.07 \\
\bottomrule
\end{tabular}
\end{center}

Young Adults contribute the highest total revenue and the highest average purchase amount. The gap is not dramatic, but it is large enough to justify closer targeting and message testing for that segment.

\findingbox
{Young Adults contribute the highest revenue and the highest average purchase amount among the four age groups.}
{This segment appears to offer the strongest near-term opportunity for targeted marketing and retention activity.}
{Allocate more campaign attention to Young Adults and tailor product selection, messaging, and channel strategy to improve conversion and retention within this segment.}

\subsection{Gender Comparison}
Male customers generated \textbf{157,890 USD} in revenue compared with \textbf{75,191 USD} from female customers. However, the average purchase amount is slightly higher for female customers:
\begin{itemize}
  \item Female average purchase: \textbf{60.25 USD}
  \item Male average purchase: \textbf{59.54 USD}
\end{itemize}

The revenue gap is driven mainly by customer count rather than by per-transaction spending power. Female customers spend slightly more per order on average, but the dataset contains many more male customer records.

\findingbox
{Male customers generate more total revenue, but female customers spend slightly more per order on average.}
{The revenue gap appears to come from customer mix rather than stronger spending per transaction.}
{Treat gender differences as a segmentation signal for messaging and acquisition analysis rather than assuming one group is inherently more valuable on an order basis.}

\subsection{Subscriptions Do Not Yet Increase Spend}
One of the most important findings is that subscribed customers do \textbf{not} currently spend more than non-subscribed customers in this dataset:
\begin{center}
\begin{tabular}{lrrr}
\toprule
\textbf{Subscription Status} & \textbf{Revenue} & \textbf{Customers} & \textbf{Avg. Spend} \\
\midrule
No & 170,436 & 2,847 & 59.87 \\
Yes & 62,645 & 1,053 & 59.49 \\
\bottomrule
\end{tabular}
\end{center}

This is an important management finding. Subscription programs are usually expected to increase loyalty and spend, but the current data does not support that assumption.

\findingbox
{Subscribers do not currently spend more than non-subscribers.}
{The current subscription program is not yet creating clear incremental customer value.}
{Redesign the subscription offer and reposition it as a retention lever rather than assuming it is already a revenue-growth lever. The business should test clearer benefits such as shipping perks, member-only discounts on high-margin products, early access to launches, or loyalty points tied to spending thresholds.}

\subsection{Repeat Buyers Are Only Moderately More Likely to Subscribe}
Customers with more than five previous purchases are somewhat more likely to subscribe, but the difference is limited:
\begin{itemize}
  \item Non-repeat buyers subscribed at roughly \textbf{22.4\%}
  \item Repeat buyers subscribed at roughly \textbf{27.6\%}
\end{itemize}

This indicates a positive relationship, but not a strong one.

\findingbox
{Repeat buyers are only moderately more likely to subscribe than non-repeat buyers.}
{The business is not fully monetizing an already engaged customer group.}
{Create more targeted conversion offers for repeat buyers, especially after a defined purchase threshold, because this group already shows higher engagement and is the most natural audience for subscription upsell.}

\subsection{Customer Segmentation}
Using previous purchase counts, customers were segmented into:
\begin{itemize}
  \item \textbf{New:} 1 previous purchase
  \item \textbf{Returning:} 2 to 10 previous purchases
  \item \textbf{Loyal:} more than 10 previous purchases
\end{itemize}

The segment distribution is:
\begin{center}
\begin{tabular}{lr}
\toprule
\textbf{Segment} & \textbf{Customers} \\
\midrule
Loyal & 3,116 \\
Returning & 701 \\
New & 83 \\
\bottomrule
\end{tabular}
\end{center}

The dataset is heavily weighted toward loyal customers, so this segmentation is useful directionally but should be interpreted carefully.

\findingbox
{The customer base is heavily concentrated in the Loyal segment, with far fewer New and Returning customers.}
{Segment insights are directionally useful, but the current data mix may overrepresent established customers.}
{Use this segmentation for targeting logic and dashboard storytelling, but validate it with future data before using it for major retention or lifecycle decisions.}

\subsection{Discount Usage}
Discounting does not simply capture low-value customers. In fact, \textbf{839} discounted purchases were still above or equal to the global average transaction value. This tells us that discounts are reaching a meaningful number of higher-spending customers as well.

The products with the highest discount usage rate are:
\begin{center}
\begin{tabular}{lr}
\toprule
\textbf{Product} & \textbf{Discount Rate} \\
\midrule
Hat & 50.00\% \\
Sneakers & 49.66\% \\
Coat & 49.07\% \\
Sweater & 48.17\% \\
Pants & 47.37\% \\
\bottomrule
\end{tabular}
\end{center}

This suggests discount strategy should be reviewed at the product level rather than treated as a uniform tactic across the catalog.

\findingbox
{Several products, including Hat, Sneakers, Coat, Sweater, and Pants, appear highly discount-dependent.}
{Some demand may be sustained by promotions rather than by underlying product strength, creating margin risk.}
{Monitor items such as Hat, Sneakers, Coat, Sweater, and Pants for margin erosion and test whether part of their demand can be preserved through improved positioning, bundling, or merchandising rather than frequent discounting.}

\subsection{Product Quality and Demand}
The top 5 products by average review rating are:
\begin{center}
\begin{tabular}{lr}
\toprule
\textbf{Product} & \textbf{Average Rating} \\
\midrule
Gloves & 3.86 \\
Sandals & 3.84 \\
Boots & 3.82 \\
Hat & 3.80 \\
Skirt & 3.78 \\
\bottomrule
\end{tabular}
\end{center}

Meanwhile, the most purchased products by category are:
\begin{center}
\begin{tabular}{lll}
\toprule
\textbf{Category} & \textbf{Top Product} & \textbf{Orders} \\
\midrule
Accessories & Jewelry & 171 \\
Clothing & Blouse / Pants & 171 each \\
Footwear & Sandals & 160 \\
Outerwear & Jacket & 163 \\
\bottomrule
\end{tabular}
\end{center}

These results help separate \textit{quality leaders} from \textit{volume leaders}, which is important because customer satisfaction and sales volume do not always overlap.

\findingbox
{High-rated products such as Gloves, Sandals, and Boots do not fully overlap with the highest-volume products.}
{Product quality signals can support conversion and trust even when they are not the top sales drivers.}
{Feature high-rated products more prominently in recommendation widgets, paid campaigns, product collections, and landing pages to capitalize on proven customer approval.}

\subsection{Shipping Type Comparison}
Express shipping users spend slightly more on average than Standard shipping users:
\begin{itemize}
  \item Express: \textbf{60.48 USD}
  \item Standard: \textbf{58.46 USD}
\end{itemize}

The gap is small, but it suggests that customers choosing faster fulfillment may be modestly more valuable.

\findingbox
{Express shipping users spend slightly more on average than Standard shipping users.}
{Faster fulfillment may be associated with a somewhat higher-value customer profile.}
{Use shipping preference as a possible premium-service or upsell signal, while validating the pattern further with additional behavioral or profitability data.}

\section{Power BI Dashboard Design}
The dashboard presents the customer behavior story in a format that non-technical stakeholders can explore quickly. The layout uses KPI cards, slicers, category-level charts, and age-based breakdowns.

\subsection{Dashboard Objectives}
The dashboard focuses on four practical goals:
\begin{itemize}
  \item give a fast overview of customer count, average purchase amount, and review rating;
  \item let users filter by subscription status, gender, category, and shipping type;
  \item compare category performance through both revenue and sales count;
  \item highlight age-group behavior to support segmentation and campaign planning.
\end{itemize}

\subsection{Business Reporting Value}
Power BI completes the workflow by turning analytical output into a decision-support layer:
\begin{itemize}
  \item Stakeholders could use filters to test hypotheses without writing code.
  \item KPI cards highlight the most important topline metrics first.
  \item Drill-down visuals support category, age, and subscription analysis.
\end{itemize}

\section{From Insight to Experimentation}
Because this dataset is transactional and does not include an experimental design, the project does not claim that A/B testing was already conducted. Instead, the analysis produces hypotheses for future testing.

\subsection{Why A/B Testing Belongs as a Next Step}
The current analysis identifies where opportunities likely exist, but not whether a specific intervention will outperform an alternative in production. A/B testing is therefore the natural next step after insight generation because it allows the business to validate decisions with measurable causal evidence.

\subsection{Recommended Experiments}
\textbf{1. Subscription Offer Test}
\begin{itemize}
  \item \textbf{Hypothesis:} Adding stronger subscriber benefits will improve conversion among repeat buyers.
  \item \textbf{Target group:} customers with more than five previous purchases.
  \item \textbf{Primary metric:} subscription conversion rate.
  \item \textbf{Secondary metrics:} average order value, repeat purchase rate.
\end{itemize}

\textbf{2. Discount Strategy Test}
\begin{itemize}
  \item \textbf{Hypothesis:} Replacing broad discounts with more selective offers on discount-heavy products can preserve demand while reducing promotion dependency.
  \item \textbf{Target group:} customers purchasing products such as Hat, Sneakers, Coat, Sweater, and Pants.
  \item \textbf{Primary metric:} conversion rate.
  \item \textbf{Secondary metrics:} average purchase amount, discount redemption rate.
\end{itemize}

\textbf{3. Young Adult Messaging Test}
\begin{itemize}
  \item \textbf{Hypothesis:} Age-targeted messaging for Young Adult customers will improve campaign response and purchase conversion.
  \item \textbf{Target group:} Young Adult segment.
  \item \textbf{Primary metric:} purchase conversion rate.
  \item \textbf{Secondary metrics:} click-through rate, revenue per customer.
\end{itemize}

\section{Portfolio Value and Skills Demonstrated}
This project demonstrates core data analyst skills:
\begin{itemize}
  \item \textbf{Analytical framing:} translating a broad business topic into measurable questions.
  \item \textbf{Data cleaning:} handling missing values, validating redundancy, and standardizing fields.
  \item \textbf{Feature engineering:} creating age-based and frequency-based analytical dimensions.
  \item \textbf{SQL problem solving:} writing queries that go beyond basic grouping, including CTEs and window functions.
  \item \textbf{Dashboard communication:} presenting findings in an accessible, business-friendly interface.
  \item \textbf{Insight generation:} moving from descriptive results to concrete recommendations.
\end{itemize}

\section{Limitations and Next Steps}
The dataset is cross-sectional, so it does not support true cohort retention analysis. The customer segmentation is heavily weighted toward loyal customers, and the absence of margin data means discount analysis focuses on revenue and demand rather than profitability.

If this project were extended, the most valuable next steps would be:
\begin{enumerate}
  \item build cohort tables for retention and repeat purchase timing;
  \item estimate customer lifetime value using purchase frequency and spend;
  \item run A/B tests on subscription offers, discount strategy, and age-targeted messaging;
  \item test clustering methods for richer behavioral segmentation;
  \item add gross margin data to evaluate promotional efficiency;
\end{enumerate}

\section{Conclusion}
This project shows how a transaction-level dataset can be converted into business decisions through Python, SQL, and Power BI. The analysis identified the main revenue-driving categories, the highest-value customer segment, underperformance in subscription, and products that may be overly dependent on discounting.

\appendix

\section{Appendix A: Key Metric Snapshot}
\begin{longtable}{p{0.58\textwidth}r}
\toprule
\textbf{Metric} & \textbf{Value} \\
\midrule
Total revenue & 233,081 USD \\
Average purchase amount & 59.76 USD \\
Average review rating & 3.75 \\
Missing review ratings before cleaning & 37 \\
Unique customers & 3,900 \\
Clothing revenue & 104,264 USD \\
Accessories revenue & 74,200 USD \\
Footwear revenue & 36,093 USD \\
Outerwear revenue & 18,524 USD \\
Young Adult revenue & 62,143 USD \\
Subscribers & 1,053 \\
Non-subscribers & 2,847 \\
High-spending discounted purchases & 839 \\
Express average order value & 60.48 USD \\
Standard average order value & 58.46 USD \\
\bottomrule
\end{longtable}

\section{Appendix B: Selected SQL Snippets}
The following examples show the style of SQL used in the project.

\subsection{Subscribers vs. Non-Subscribers}
\begin{lstlisting}[language=SQL]
SELECT subscription_status,
       COUNT(customer_id) AS total_customers,
       ROUND(AVG(purchase_amount_usd), 2) AS avg_spend,
       ROUND(SUM(purchase_amount_usd), 2) AS total_revenue
FROM customer
GROUP BY subscription_status
ORDER BY total_revenue DESC, avg_spend DESC;
\end{lstlisting}

\subsection{Customer Segmentation}
\begin{lstlisting}[language=SQL]
WITH customer_type AS (
    SELECT customer_id,
           previous_purchases,
           CASE
               WHEN previous_purchases = 1 THEN 'New'
               WHEN previous_purchases BETWEEN 2 AND 10 THEN 'Returning'
               ELSE 'Loyal'
           END AS customer_segment
    FROM customer
)
SELECT customer_segment,
       COUNT(*) AS number_of_customers
FROM customer_type
GROUP BY customer_segment;
\end{lstlisting}

\subsection{Top 3 Products Within Each Category}
\begin{lstlisting}[language=SQL]
WITH item_counts AS (
    SELECT category,
           item_purchased,
           COUNT(customer_id) AS total_orders,
           ROW_NUMBER() OVER (
               PARTITION BY category
               ORDER BY COUNT(customer_id) DESC
           ) AS item_rank
    FROM customer
    GROUP BY category, item_purchased
)
SELECT item_rank, category, item_purchased, total_orders
FROM item_counts
WHERE item_rank <= 3;
\end{lstlisting}

\end{document}
